\section{Discussion}\label{sec:discussion}

The model isolates a specific coordination problem that can arise when the geography or organization of downstream activity changes faster than the location of information production. Its usefulness depends on keeping that mechanism separate from other forces that often accompany channel change.

\subsection{What cross-route reuse means}

Cross-route reuse does not require that raw information be literally transferred to route $R$. The maintained primitive is that information generated at $S$ improves an upstream product or design and that the improved product is then valuable through $R$. The route-$R$ term in condition~\eqref{eq:G} therefore captures the marginal system value of information after upstream reuse. This distinction matters empirically: one may observe product, diagnostic, software, quality, or design improvements rather than a direct flow of information from $S$ to $R$.

The mechanism is strongest when the expanding route can use improvements generated from information at the shrinking node without compensating that node for the resulting value. It is weakest when information is highly route-specific. The limiting case $D_R'=0$ removes the positive coordinated response altogether.

\subsection{Organization and intervention}

The analysis does not imply that a vertically coordinated organization or a policy intervention is always desirable. Coordination raises effort relative to the decentralized information producer in the welfare benchmark, but the social optimum remains higher because coordinated operating profits do not incorporate all consumer-surplus gains. Conversely, retaining a high-cost downstream node purely to preserve information production can itself be inefficient away from the local threshold region described in Section~\ref{sec:welfare}. The model therefore identifies a wedge, not a universal organizational prescription.

A practical implication is that organizational responses can target the information-production margin rather than the transaction route itself. Contracting for diagnostic effort, funding reporting infrastructure, preserving engineering interfaces, or separating information-production compensation from downstream sales are possible responses in richer environments. These are interpretations, not results of the present model, because effort is deliberately noncontractible here and no explicit contracting stage is solved.

\subsection{Empirical content}

The frozen predictions in Section~\ref{sec:institutional} suggest a way to distinguish the mechanism from a generic decline of a displaced route. Evidence consistent with the theory would combine a falling local information-production incentive with rising upstream or system demand for reusable information, and the divergence should be stronger where information is more portable across routes. Finding that both private and coordinated information incentives decline when activity moves away from $S$ would instead be consistent with weak or absent cross-route reuse.

A full empirical test would require measuring information-producing activity separately from downstream sales and addressing the endogeneity of route reallocation. Those identification problems are outside the present theory paper. The model's role is to define the comparative-static objects and the conditions under which the sign reversal should or should not appear.

\subsection{Deliberate abstractions}

Several features are held outside the model. Prices are not strategic in the core theorem. Downstream route shares are exogenous rather than chosen by firms or consumers. There is no bargaining over information, no uncertainty or signal-revelation game, and no dynamic stock of accumulated knowledge. These margins can matter in particular applications, but adding them would change the strategic architecture and is unnecessary for the result established here.

The main conclusion should therefore be read conditionally. Downstream reallocation can create a larger information-appropriability problem when the information producer's footprint shrinks and reusable information has sufficiently high marginal value through the expanding route. The model does not say that every form of channel substitution has this effect.
